\newpage
\cleardoublepage
\chapter{Discussion and Conclusion}\label{chapter:discussion_conclusion}

\Cref{chapter:results} reported the measurements. This chapter reads them back against the four research questions of \Cref{sec:objective}, fixes what the evidence cannot carry, and proposes the work that follows from those limits. \Cref{sec:discussion} restates each question verbatim and answers it from the results chapter alone, without introducing new data. \Cref{sec:limitations} collects the boundaries and the one protocol deviation that bound every answer. \Cref{sec:further_work} maps each limitation onto a concrete follow-on, and \Cref{sec:conclusion} closes with what this work establishes.

%%%%%%%%%%%%%%%%%%%%%%%%%%%%%%%%%%%%%%%%%%%%%%%%%
\section{Discussion}\label{sec:discussion}

\Cref{sec:objective} posed four research questions. They are restated here verbatim and answered one at a time.

\begin{enumerate}
    \item Does knowledge distillation from a fine-tuned teacher model into a nano-scale student model yield higher detection accuracy than fine-tuning the same student directly on the wildlife dataset, given the domain shift from COCO-style classes to $225$ mammal species?
    \item Which synthetic-image generator produces the most useful training data for this domain, and do the two criteria available before any detector is trained, apparent visual realism and per-image price, predict that ranking?
    \item How does the composition of a class's training data, from synthetic images alone through a large pool of real photographs, affect detection accuracy, and does accuracy measured on the mixed test set hold up on the real test set alone?
    \item Does a nano-scale detector meet the per-frame latency budget of the target hardware without quantization or any other compression step, measured on a target-adjacent proxy device?
\end{enumerate}

\paragraph{Question $1$: Distillation Against Direct Fine-Tuning}
No, at the one design point that was tested. \Cref{sec:results_kd_basic_comparison} compares Setup A, which blends the frozen fine-tuned teacher's cached class distribution into the student's classification target at temperature $T = 4$ and weight $\alpha = 0.5$, against Setup B, which takes the assigner's own targets untouched. Direct fine-tuning is ahead in all $17$ reported cells of \Cref{tab:kd-vs-direct}. On the mixed test set it reaches a mAP of $0.599$ against the distilled student's $0.560$, and on the real test set $0.529$ against $0.479$. The margin is widest on Band A, whose supervision is entirely synthetic and thin, at $0.086$ mixed and $0.061$ real. That is the band the distillation hypothesis was expected to help most, so the negative answer holds on its own best case rather than only on aggregate.

Three properties bound how far that answer travels, and the first of them is a confound rather than a scope limit. The distilled run trained at an effective batch size of $16$ against $32$ for the direct run, both auto-reduced from a nominal $64$ by the memory probe, so batch size is a second varied factor whose contribution to the gap is unmeasured. The distilled run also ended early, at epoch $186$ of $200$ after a tracking-client exception, with its validation metric still rising from $0.724$ at epoch $175$ to $0.733$ at epoch $185$. The direct run stopped on its own patience criterion and was the more nearly converged of the two. Neither run carries a confidence interval, so the direction is treated as real and the magnitude as unquantified.

The scope limits are narrower still. One $(T, \alpha)$ point ran rather than the grid of \Cref{sec:kd_loss}, one teacher and student pairing was tested at a parameter ratio of roughly $72$ to $1$, and only response-based distillation was used. That ratio sits far outside the range the capacity-gap literature treats as safe, and it is a documented setting for negative transfer \cite{mirzadehImprovedKnowledgeDistillation2020}, so an unfavorable outcome there is closer to the expected case than to a refutation of the mechanism. The domain shift the question names as its premise was present by construction, since the student started from COCO weights and finished on $225$ mammal species, but its size is not measured anywhere in this work. The student's zero-shot baseline was never produced (\Cref{sec:results_general_observations}), so both setups are reported as absolute levels rather than as gains over the floor they shared. What the comparison rules out is a benefit from response-based distillation at $T = 4$ and $\alpha = 0.5$ with this pairing, and not knowledge distillation for this domain.

\paragraph{Question $2$: Which Generator, and Whether Realism or Price Predicts It}
\texttt{gpt-image-2} at its cheapest quality tier, conditioned on the unabridged \texttt{full} prompt regime, produced the most useful training data of the twelve cells compared, at $0.134 \pm 0.014$ mAP on the real test set (\Cref{tab:generator-ranking}). Neither criterion the question names predicts that ranking.

Apparent visual realism selected a reasonable generator that was not the best available one. The incumbent was chosen for the production synthetic dataset because its output looked more realistic, and it ranks third of twelve. \texttt{gpt-image-2} at its cheapest tier beats it by roughly $28\%$ under an identical prompt regime, and its own vendor's cheaper variant is statistically tied with it. Per-image price predicts nothing at all. The two \texttt{gpt-image-2} quality tiers differ by a factor of $2.9$ in price and are separated by $0.008$ mAP against a combined standard deviation of roughly $0.018$ (\Cref{sec:results_cost_vs_quality}). In the local grid the most expensive cell to produce, at nearly $48$ GPU-hours, ranks eleventh of twelve, while the cheapest by two orders of magnitude ranks twelfth. In neither vendor family is the cheaper tier the worse one, so generation cost is not usable as a proxy for training-data quality.

A third criterion the question did not anticipate outranks both. Compressing the prompt costs roughly half the downstream accuracy, at $-45\%$ and $-48\%$ for the two models of the controlled ablation in \Cref{tab:prompt-length-ablation}, and within-group confusion over the zebras rises in step. That effect is larger than most model-to-model differences in \Cref{tab:generator-ranking} and moves the top-ranked cell to fifth place. How much diagnostic detail a generator is given therefore matters more than which generator is chosen, across the range of models tested.

The ranking answers the question for data-rich classes and under the single evaluation axis that was executed. On the rare classes it answers nothing. Five of the six stay below $0.05$ average precision for every cell in the grid (\Cref{sec:results_per_class_structure}), and one generator draws kinkajous with a ringed tail and a spotted coat at no measured penalty (\Cref{fig:kinkajou-generator-samples}). The honest form of the answer is that the comparison ranks generators for well-photographed classes and establishes that its own instrument cannot rank them for the long-tail classes that motivated generating images in the first place.

\paragraph{Question $3$: Training-Data Composition, and Whether Mixed Accuracy Holds Up}
Composition matters, and it matters non-linearly. Bands A, B and C share an equal $200$-image training budget and differ only in what those images are, and \Cref{tab:band-grid-real} reports them on the real test set. Replacing half of a small real training set with synthetic imagery costs nothing measurable in two of the three models. \textit{YOLO26n} under direct fine-tuning reaches $0.588$ under Band B against $0.591$ under Band C at \texttt{fine} granularity, and the fine-tuned teacher ensemble $0.638$ against $0.639$. \textit{YOLOv5s} is the exception at $0.341$ against $0.386$, favoring real photographs by $0.045$. At \texttt{coarse} granularity Band B sits above Band C by $0.019$ for the primary detector and by $0.034$ for the ensemble.

Replacing all of it is a different matter. Band A leaves the primary detector at $0.300$ on the real test set, roughly half the $0.591$ the same budget of real photographs achieves, and \textit{YOLOv5s} at $0.058$. That gap is an order of magnitude larger than the Band B against Band C difference beside it, so its direction survives the absence of confidence intervals where the smaller difference does not. Two qualifications apply to both readings. Each figure comes from a single training run, and the bands hold different species assigned by data availability rather than at random (\Cref{sec:band_structure}), so every cross-band difference is confounded with how hard the species in each band happen to be. Whether Band A fails on image realism or on the diversity of poses and backgrounds a generator produces is not separable from these numbers either.

The second half of the question has a blunter answer. Accuracy measured on the mixed test set does not hold up on the real test set alone. The real-only breakout falls below the mixed headline for all four models, by between $0.042$ and $0.089$ mAP (\Cref{tab:headline-cross-model}). The paired per-class delta of \Cref{tab:domain-shift-delta} shows why. For the primary detector $\Delta_{\text{domain}}$ is negative for $208$ of the $224$ banded classes, and for all $76$ classes of Band A and Band B without exception. The synthetic test set is the easier domain by a median of $0.21$ mAP even on data-rich Band D, rising to $0.72$ on Band A, and the band ordering is identical across all three fine-tuned models.

Roughly $0.2$ mAP of that gap is not domain shift. A synthetic test image holds one centered animal by construction while a real one averages $1.44$ boxes (\Cref{sec:results_regime_comparison}), and the off-the-shelf ensemble, which never saw a synthetic image in training, still scores higher on the synthetic domain in every band. Only the band-dependent remainder above that offset reads as an effect of how a class was trained, on the assumption that the difficulty offset carries over to the trained models.

The watchdog that \Cref{sec:mixed_default} attached to the mixed default therefore fired. Its verdict is recorded in \Cref{sec:results_domain_shift}: the discrepancy is large, systematic and one-directional, and it inflates the mixed headline most for exactly the classes the mixed default was introduced to stabilize. The designed revision, which is to restore the real test set as the reported default and retain the mixed figure as secondary, was not carried out, because rewriting every figure in \Cref{chapter:results} did not fit the available time. That is a deviation from the stated evaluation protocol rather than an unexplored axis, and \Cref{sec:limitations} treats it as one. Two properties limit the damage without repairing it. Every headline figure in that chapter appears with its real-only breakout beside it, so no claim there rests on the mixed figure alone. And switching the default moves the cross-model ranking in one place only, where \textit{YOLOv5s} and the off-the-shelf ensemble exchange positions.

\paragraph{Question $4$: The Per-Frame Latency Budget Without Compression}
No. No configuration measured in \Cref{sec:results_embedded_benchmark} meets the project's $\leq 30$\,ms per-frame design target, and the shortfall is not marginal. The fastest model is \textit{YOLO26n} under direct fine-tuning, at $423$\,ms end-to-end on the \textit{Raspberry Pi 400} proxy hardware, which the proxy relationship of \Cref{sec:lit_embedded_inference} translates to roughly $360$--$381$\,ms on the \textit{QCS605} CPU. That is about twelve times the budget. \textit{YOLOv5s}, the architecture family closest to the detector already shipping in the production binocular, needs $949$\,ms. The memory axis is met, with all three detectors below the $\leq 500$\,MB ceiling at $284$\,MB for either \textit{YOLO26n} run and $362$\,MB for \textit{YOLOv5s}, while neither teacher component is.

The teacher's cost settles a separate question along the way. \textit{MegaDetector v5a} alone takes $31.0$\,s per frame and peaks at $1{,}570$\,MB, and the composed ensemble estimate comes to roughly $35$\,s per frame, three orders of magnitude above the budget. \Cref{sec:lit_yolo_family}'s reasoned argument that the ensemble cannot ship on this class of hardware becomes a measured one.

The question's own wording carries the reason this answer is not a verdict on the hardware. It asks about a detector without quantization or any other compression step, and that omission is a scope decision recorded in \Cref{sec:results_kd_and_embedded} rather than an attempt that failed. The figures are full-precision and unoptimized, so they bound latency from above rather than estimating what the target hardware can do. The \textit{QCS605}'s Hexagon 685 DSP has no analogue on the proxy hardware, and the accelerated path a shipping product would use, for which $8$-bit quantization is reported at up to $10\times$ over floating-point CPU inference on this class of DSP \cite{krishnamoorthiQuantizingDeepConvolutional2018}, remains entirely unmeasured. Whether $30$\,ms is reachable on the target is therefore still open.

\paragraph{Summary}
Read together, the four answers point one way. The two questions that received the bulk of the experimental effort returned usable findings, one negative about the criteria a practitioner would reach for first when choosing a generator, and one positive about a specific and bounded use of synthetic imagery. The two questions answered at a single design point returned clean negatives inside that scope, for distillation and for full-precision deployment alike. None of the four answers describes an accuracy problem for the chosen student architecture, and the one that blocks deployment is a compression problem.

%%%%%%%%%%%%%%%%%%%%%%%%%%%%%%%%%%%%%%%%%%%%%%%%%
\section{Limitations}\label{sec:limitations}

\Cref{sec:results_general_observations} fixes five scope boundaries and \Cref{sec:results_domain_shift} records one protocol deviation. This section collects those, adds the qualifications that cut across chapters, and states for each what it costs the corresponding finding.

\paragraph{No confidence intervals anywhere}
The bootstrap resampling specified in \Cref{sec:metric_selection} was never run, since the resample count stayed at zero in every evaluation run, and every model contributes a single training run rather than a repeated trial. No difference reported anywhere in \Cref{chapter:results} can therefore be called statistically significant. The consequence is specific rather than uniform. A direction holding across $17$ partly independent views of one pair of runs, or a gap of the size separating Band A from Band C, is unlikely to be seed noise. A difference of $0.003$ between Band B and Band C carries no information about which composition is better, which is why that result is stated as a cost of nothing measurable rather than as a tie or a win. The seed episode of \Cref{sec:results_practicability_verdict} shows a single-seed gap in exactly that size range reversing outright once two further seeds were added.

\paragraph{The watchdog revision was not carried out}
This is the one item here that deviates from a stated protocol rather than leaving an axis unexplored. \Cref{sec:mixed_default} committed in advance to revising the default evaluation axes if a systematic discrepancy between the mixed and real results appeared. The discrepancy appeared, it is large and one-directional, and the revision did not follow. Every mixed-set figure in this work should accordingly be read as a ranking instrument and not as an estimate of accuracy on real photographs, with the real-only breakout beside it carrying that second reading.

\paragraph{No student zero-shot baseline}
An evaluation of \textit{YOLO26n} from raw COCO weights failed on a missing data directory and was not retried after the directory was repaired. Every student figure in this work is therefore an absolute level and not a gain over an untuned starting point. What fine-tuning on this domain bought is unmeasured, which bears directly on Question $1$: Setup A and Setup B are compared to each other and neither is anchored to the floor they both started from.

\paragraph{No distillation hyperparameter grid, one pairing, one distillation type}
Only the default $(T, \alpha)$ point ran, so \Cref{tab:kd-vs-direct} is one cell of the grid \Cref{sec:kd_loss} describes. One teacher and one student were paired, and feature-based alignment was excluded on architectural grounds before any run started, so only response-based distillation was tested. Reading Question $1$'s negative answer as a result about knowledge distillation for fine-grained wildlife detection would overstate it by the size of the space \Cref{sec:lit_knowledge_distillation} surveys, including the attention-aware detection variants that literature reports gains for \cite{yangFocalGlobalKnowledge2022}.

\paragraph{The multi-animal subset was never evaluated}
An image holding several animals is where a teacher's per-image class distribution carries the most information the assigner's own hard targets do not, so that cut is the concrete test of distillation's expected supervision advantage. It is also the setting the per-image teacher cache of \Cref{sec:kd_loss} approximates most crudely. Its absence means the one mechanism that could have favored Setup A was not probed, which is a further reason to read Question $1$'s answer as bounded.

\paragraph{No quantization-aware training}
No quantization, pruning or compression pass was applied, and \Cref{sec:results_kd_and_embedded} records that as a scope decision rather than an attempted step. The full-precision figures of \Cref{tab:embedded-latency} bound latency from above and do not estimate what the target hardware can do, particularly since the accelerated path on the \textit{QCS605} is a quantized one \cite{jacobQuantizationTrainingNeural2018}. Question $4$'s answer is a feasibility statement about unoptimized models on a target-adjacent CPU, and calling it a deployment verdict would require that pass.

\paragraph{Every teacher figure is an upper bound}
The ground-truth boxes of both test domains were produced by \textit{MegaDetector} (\Cref{sec:gt_annotation_integrity}), and the teacher ensemble localizes with that same detector, so its localization is scored partly against its own output. Its mAP averaged over intersection-over-union thresholds from $0.5$ to $0.95$ equals its mAP50 to three decimal places, which no single-stage detector reported here approaches. Every teacher figure in \Cref{chapter:results} is therefore an upper bound, and the measured teacher-to-student gap is an upper bound with it. A second caveat applies to the teacher's fine-tuning gain of $0.176$ mixed and $0.154$ real, since part of the fine-tuning corpus was admitted by agreeing with the pre-fine-tuning classifier (\Cref{sec:teacher_finetuning}). That gain is optimistic by an unknown and probably small amount.

\paragraph{Six qualifications on the generator comparison}
Beyond the deferred evaluation axes, \Cref{sec:results_practicability_verdict} states six limits on the sub-study. The annotations are automatic and un-reviewed throughout. Seed counts are uneven, at three per API cell against two per local cell, short of the three the protocol itself recommends. The design is not a full factorial, since the local tier has no unabridged-prompt cell and the API tier no capacity-matched one. Evaluation is real-only, with no mixed condition defined for the sub-study. One planned cell, the quality ceiling from the more expensive tier of the incumbent's vendor, was never generated, so the grid has no upper anchor. And three fairness asymmetries remain, namely a prompt trimmed for one vendor's character cap on long-article species, negative prompts unavailable to one local cell, and two prompt regimes that are not seed-matched to each other. Together these mean \Cref{tab:generator-ranking} supports the ordinal claims of Question $2$ and supports neither an absolute accuracy figure nor a full crossing of generator against prompt length.

\paragraph{The blind rating rubric was never run}
The multi-rater rubric specified in \Cref{sec:generator_comparison_evaluation} was not carried out, nor were three of the four automatic proxies, and the one proxy computed discriminates too weakly to separate the local cells. The cost is precise rather than diffuse. The rubric was designed as the instrument for the regime in which downstream accuracy has no resolution, which is the rare classes, so its absence is why a visually obvious species error can be shown directly in \Cref{fig:kinkajou-generator-samples} and detected by no executed measure.

\paragraph{Error types are not separated, and training cost is not comparable}
The error-type decomposition named in \Cref{sec:metric_selection} was left as a future add-on rather than integrated into the evaluation suite, so cross-group misclassification, duplicate detections and background false positives all land in $\Delta_{\text{coarse}}$ together (\Cref{sec:results_granularity_gaps}). $\Delta_{\text{coarse}}$ is an upper bound on cross-group species confusion and not a measurement of it. The finding that classification rather than localization limits accuracy survives this, because the recall figures of \Cref{tab:headline-coco12} support it independently, while no claim about which kind of classification error dominates does. Separately, the wall-clock cost of the four final runs is reported in \Cref{sec:results_general_observations} and is not a throughput comparison, since every run overlapped on the same device with at least one other.

%%%%%%%%%%%%%%%%%%%%%%%%%%%%%%%%%%%%%%%%%%%%%%%%%
\section{Further Work}\label{sec:further_work}

Each item below follows from one of the limitations above, and the first four bear on the deployment question that is currently unresolved.

\paragraph{Quantization and the accelerated path}
The largest remaining lever on the embedded target is a compression pass, since the accelerated path on the \textit{QCS605} is a quantized one and the full-precision figures of \Cref{tab:embedded-latency} miss the budget by roughly twelve times. Post-training quantization is the cheaper first step, reported to hold classification accuracy within $2\%$ of the floating-point baseline under per-channel weight and per-layer activation $8$-bit quantization \cite{krishnamoorthiQuantizingDeepConvolutional2018}, and quantization-aware training narrows the residual gap further by adapting weights to quantization noise during training \cite{jacobQuantizationTrainingNeural2018}. Both map onto the vendor toolchain named in \Cref{sec:lit_embedded_inference} \cite{QualcommAimet2026}. The measurement that settles Question $4$ has to be taken on the Hexagon 685 DSP itself, because no proxy device used in this work has an analogue to it.

\paragraph{Input resolution as a second latency lever}
Reducing the input resolution the detector runs at is an independent direction for closing the remaining latency gap, and it is named here as a direction rather than reported as a result. A convolutional detector's forward cost scales with the pixel count it is given, and the pre-processing measurements of \Cref{sec:results_embedded_benchmark} already show decode and letterboxing costs scaling with capture resolution in the same way. Any such reduction needs a resolution-native training arm beside it, since a detector trained at one input size and then run at a smaller one loses accuracy for reasons unrelated to the resolution choice itself. No figure from this direction is reported in \Cref{chapter:results}.

\paragraph{The distillation grid and other pairings}
Question $1$'s negative answer is scoped to one cell, so the direct follow-on is the sweep over temperature $T$ and blend weight $\alpha$ that \Cref{sec:kd_loss} was written around. Two further axes matter more than the sweep. Closing the parameter ratio of roughly $72$ to $1$ with an intermediate-capacity teacher assistant is the remedy the capacity-gap literature proposes for this exact configuration \cite{mirzadehImprovedKnowledgeDistillation2020}, and a smaller fine-tuned detector used as the teacher would test the same idea more cheaply. Attention-aware detection distillation, which separates foreground from background rather than matching outputs uniformly, is the other \cite{yangFocalGlobalKnowledge2022}. Any repeat should hold the effective batch size equal across the two setups, which the comparison reported here did not.

\paragraph{The student's zero-shot floor}
Evaluating \textit{YOLO26n} from raw COCO weights on the same test set is one evaluation run and the cheapest missing measurement in this work. It would turn every student figure from an absolute level into a gain, and it would quantify the domain shift Question $1$ names as its premise instead of assuming it.

\paragraph{The multi-animal subset}
Cutting the test set to images holding more than one animal and re-scoring Setup A against Setup B on that subset alone is the concrete test of distillation's expected supervision advantage. Both runs and both prediction sets are archived, so the cut needs no retraining, and a real test image averages $1.44$ boxes, so the subset is not a marginal slice. It also tests the one implementation shortcut \Cref{sec:kd_loss} took, since a per-image teacher cache approximates a multi-animal image most crudely.

\paragraph{The blind rating rubric and the deferred proxies}
The rubric of \Cref{sec:generator_comparison_evaluation} remains fully specified and directly executable on the archived images, and it needs raters rather than compute. It is the axis that would resolve the rare-class question the executed axis has no power over, since a generator drawing the wrong animal is obvious to a human rater and invisible to a $100$-image fine-tune scored on a few dozen real photographs. Teacher recognition confidence, per-class distributional distances and prompt-adherence scoring are executable on the same archive and were specified alongside it.

\paragraph{Confidence intervals, computed retroactively}
Bootstrap resampling over the stored per-image predictions would give every comparison in this work the significance test it currently lacks, without retraining anything. The evaluation suite already accepts a resample count and the prediction artifacts are archived, so the cost is compute rather than redesign. The comparisons that need it most are the ones whose gaps are small, namely Band B against Band C and the two statistically tied generator tiers. Repeating the single-run models at further seeds is the more expensive and more complete version of the same fix.

\paragraph{Carrying out the watchdog revision}
The revision \Cref{sec:mixed_default} committed to remains outstanding, which is to make the real test set the reported default across every table in \Cref{chapter:results} and retain the mixed figure as secondary. Nothing blocks it but the rewriting effort, since every table already holds both domains, and doing it would move the cross-model ranking in one place only. Redesigning the mixed condition is the larger version of the same task, because the current synthetic test set holds one centered animal per image and is easier for that reason alone rather than only through domain shift.

%%%%%%%%%%%%%%%%%%%%%%%%%%%%%%%%%%%%%%%%%%%%%%%%%
\section{Conclusion}\label{sec:conclusion}

This work set out to build, train and evaluate a nano-scale wildlife species detector for embedded deployment, over a class universe of $225$ mammal species of which $50$ retain fewer than $150$ usable real photographs each. Four research questions organized it, and the effort behind them was deliberately uneven. Roughly four fifths of the available time went into dataset construction and synthetic-image generation, so the substantive contribution sits in the two sub-studies that effort produced, while the distillation and deployment questions were each answered at a single design point.

The first sub-study is a controlled comparison of twelve generator and prompt-regime cells, each supplying $1{,}200$ synthetic training images of the same $12$ classes to the same detector under the same schedule. Neither criterion available before a detector is trained predicts what the resulting detector achieves. The generator chosen for the production dataset on apparent realism ranks third of twelve, and per-image price is unrelated to outcome in both vendor families and in the local grid. Prompt detail outranks generator identity, since compressing the prompt costs roughly half the downstream accuracy while the spread across generators under their best prompt is smaller. For the long-tail classes that motivated generating images at all, the comparison establishes that its own instrument has no resolution.

The second is the band comparison at a matched $200$-image training budget. Replacing half of a small real training set with synthetic imagery costs nothing measurable for the primary detector or for the fine-tuned teacher, at $0.588$ against $0.591$ mAP on the real test set, and costs $0.045$ for \textit{YOLOv5s}. Replacing all of it costs roughly half the accuracy, at $0.300$ against $0.591$. Both figures rest on single runs and on bands that hold different species, so they are directional rather than exact. The paired per-class domain-shift delta then fired the watchdog attached to the mixed default, and the designed revision did not follow within the available time, which is the one protocol deviation this work discloses.

On the central question, distillation from the fine-tuned teacher lost to direct fine-tuning in all $17$ reported cells, at $0.599$ against $0.560$ mAP on the mixed test set and $0.529$ against $0.479$ on the real test set. That is a negative result at one $(T, \alpha)$ point, with one teacher and student pairing at a parameter ratio of roughly $72$ to $1$, and three disclosed asymmetries keep it from being a clean single-factor comparison. It is nonetheless a useful result for anyone weighing the same architecture-scale gap, because the cheaper recipe won on every axis measured and won by the widest margin on the thin-data band where the teacher's soft targets were expected to help most.

The practical status is easy to state. \textit{YOLO26n} under direct fine-tuning is at once the smallest, the most accurate and the fastest detector measured here, at $2.71\,\mathrm{M}$ parameters, $0.599$ mAP on the mixed test set with $0.529$ on the real test set, and $423$\,ms per frame on the proxy hardware. That last figure is roughly twelve times the target hardware's per-frame budget, and no full-precision model measured in this work comes closer. Nothing about this detector's accuracy stands between the work reported here and a shippable model. A compression pass does, and the accelerated path such a pass would target has not been measured at all.
